% Options for packages loaded elsewhere
\PassOptionsToPackage{unicode}{hyperref}
\PassOptionsToPackage{hyphens}{url}
\documentclass[
  9pt,
]{ltjsarticle}
\usepackage{xcolor}
\usepackage[margin=14mm]{geometry}
\usepackage{amsmath,amssymb}
\setcounter{secnumdepth}{-\maxdimen} % remove section numbering
\usepackage{iftex}
\ifPDFTeX
  \usepackage[T1]{fontenc}
  \usepackage[utf8]{inputenc}
  \usepackage{textcomp} % provide euro and other symbols
\else % if luatex or xetex
  \usepackage{unicode-math} % this also loads fontspec
  \defaultfontfeatures{Scale=MatchLowercase}
  \defaultfontfeatures[\rmfamily]{Ligatures=TeX,Scale=1}
\fi
\usepackage{lmodern}
\ifPDFTeX\else
  % xetex/luatex font selection
\fi
% Use upquote if available, for straight quotes in verbatim environments
\IfFileExists{upquote.sty}{\usepackage{upquote}}{}
\IfFileExists{microtype.sty}{% use microtype if available
  \usepackage[]{microtype}
  \UseMicrotypeSet[protrusion]{basicmath} % disable protrusion for tt fonts
}{}
\makeatletter
\@ifundefined{KOMAClassName}{% if non-KOMA class
  \IfFileExists{parskip.sty}{%
    \usepackage{parskip}
  }{% else
    \setlength{\parindent}{0pt}
    \setlength{\parskip}{6pt plus 2pt minus 1pt}}
}{% if KOMA class
  \KOMAoptions{parskip=half}}
\makeatother
\setlength{\emergencystretch}{3em} % prevent overfull lines
\providecommand{\tightlist}{%
  \setlength{\itemsep}{0pt}\setlength{\parskip}{0pt}}
\setlength{\parskip}{2pt}\setlength{\parindent}{0pt}
\usepackage{bookmark}
\IfFileExists{xurl.sty}{\usepackage{xurl}}{} % add URL line breaks if available
\urlstyle{same}
\hypersetup{
  hidelinks,
  pdfcreator={LaTeX via pandoc}}

\author{}
\date{}

\begin{document}

\section{中心投影・球面投影の幾何 ──
多軸モデルと、その物理的舞台への一瞥（観察ノート）}\label{ux4e2dux5fc3ux6295ux5f71ux7403ux9762ux6295ux5f71ux306eux5e7eux4f55-ux591aux8ef8ux30e2ux30c7ux30ebux3068ux305dux306eux7269ux7406ux7684ux821eux53f0ux3078ux306eux4e00ux77a5ux89b3ux5bdfux30ceux30fcux30c8}

\textbf{木原 範昭}（Noriaki Kihara）／ WF System Co., Ltd.／ 大阪大学
基礎工学部 卒 ORCID 0009-0004-6753-4020 ／ CC BY 4.0 ／ 2026-06

\begin{quote}
\textbf{観察ノート。} 新規の幾何学的定理は主張しない。古典的写像（radial
projection）を整理し、その上に現れる構造を観察する。物理的導出は行わない。
\end{quote}

\subsection{\texorpdfstring{1. 球面投影 \(\sigma_R\)（radial
projection）}{1. 球面投影 \textbackslash sigma\_R（radial projection）}}\label{ux7403ux9762ux6295ux5f71-sigma_rradial-projection}

\[\sigma_R:\mathbb{R}^{n+1}\setminus\{0\}\to S^n(R),\qquad \sigma_R(x)=\frac{R}{\|x\|}\,x\]

\begin{itemize}
\tightlist
\item
  位相幾何の \textbf{radial projection}（Hatcher \emph{Algebraic
  Topology} Ch.0 の変形レトラクト）と同一。\(C^\infty\)
  レトラクションで冪等 \(\sigma_R\circ\sigma_R=\sigma_R\)。
\item
  商空間
  \((\mathbb{R}^{n+1}\setminus\{0\})/\mathbb{R}_{>0}\cong S^n(R)\)。微分の核
  \(=\) 動径方向 \(\mathrm{span}\{x\}\)、像
  \(=x^{\perp}=T_{\sigma_R(x)}S^n(R)\)（階数 \(n\)
  の沈め込み）。角度保存、半径スケールに対する方向不変。
\end{itemize}

\subsection{\texorpdfstring{2. 中心投影 \(\Phi_R\)（gnomonic ＝
球面投影の接平面への制限）}{2. 中心投影 \textbackslash Phi\_R（gnomonic ＝ 球面投影の接平面への制限）}}\label{ux4e2dux5fc3ux6295ux5f71-phi_rgnomonic-ux7403ux9762ux6295ux5f71ux306eux63a5ux5e73ux9762ux3078ux306eux5236ux9650}

北極で接する接超平面 \(\Pi_R=\{x_{n+1}=R\}\cong\mathbb{R}^n\) への制限：

\[\Phi_R=\sigma_R\big|_{\Pi_R}:\ \Pi_R\ \xrightarrow{\ \sim\ }\ S^n_+(R)\quad(\text{開上半球面・微分同相})\]

\begin{itemize}
\tightlist
\item
  引き戻し計量
  \(g_{\mu\nu}=\dfrac{R^2}{\ell^2}\!\left(\delta_{\mu\nu}-\dfrac{x_\mu x_\nu}{\ell^2}\right)\)（\(\ell=\sqrt{R^2+|x|^2}\)）は
  \(G_{\mu\nu}+\Lambda g_{\mu\nu}=0\)（\(\Lambda=\tfrac{(n-1)(n-2)}{2R^2}\)）を満たし、\textbf{de
  Sitter 空間の Beltrami 座標と内在的に一致}。\(R\to\infty\)
  で平坦（Minkowski）に退化。
\item
  \textbf{核心}：\(\sigma_R\)（非単射・全球）と
  \(\Phi_R\)（単射・開上半球のみ）の対比。
\end{itemize}

\subsection{3. 多軸モデル}\label{ux591aux8ef8ux30e2ux30c7ux30eb}

背景 \(\mathbb{R}^{n+1}\) の \(n+1\)
本の軸のどれを投影中心に選んでも、計量・曲率テンソルの構造は同一（\textbf{軸の対等性}）。\(n+1\)
通りの主観空間が同時に立ち、相互変換は
\(T_{A\to B}=\Phi_B^{-1}\circ\Phi_A\)。観測者の立場により「投影中心軸（直接アクセス不可）」と「主観座標軸」の役割が交換される。

\subsection{4.
物理的舞台への一瞥（観察のみ）}\label{ux7269ux7406ux7684ux821eux53f0ux3078ux306eux4e00ux77a5ux89b3ux5bdfux306eux307f}

外部を持たない真空（内在幾何系）の中心投影として、等質な4次元球面
\(S^4(R_{\mathrm U})\)（断面曲率 \(1/R_{\mathrm U}^2\)、\(dS_4\)
のユークリッド版）を取ると、その上に「広がり位相を持つ局所モード（粒子様状態）」\(P_a=(\boldsymbol{\nu}_a,R_a)\)、広がり
\(W_a=2R_a\)
を置けることが観察される。ローレンツ計量・因果構造・物理的導出は主張しない。［Paper
14］

\subsection{参考（Concept DOI）}\label{ux53c2ux8003concept-doi}

\begin{itemize}
\tightlist
\item
  球面投影 \(\sigma_R\)
  の定義（基礎写像）：\href{https://doi.org/10.5281/zenodo.20462569}{10.5281/zenodo.20462569}
\item
  中心投影による4次元空間の幾何学的定式化：\href{https://doi.org/10.5281/zenodo.19427780}{10.5281/zenodo.19427780}
\item
  中心投影の合成演算（多軸・合成曲率半径の閉形式）：\href{https://doi.org/10.5281/zenodo.20060728}{10.5281/zenodo.20060728}
\item
  ［Paper
  14］真空宇宙の中心投影と、広がり位相を持つ粒子様状態：\href{https://doi.org/10.5281/zenodo.20543044}{10.5281/zenodo.20543044}
\end{itemize}

GitHub: github.com/WurabeSeiji/ai-chat-logs-open ／ 全文は Zenodo
に公開（日英対訳・CC BY 4.0）

\end{document}
